\scanpage{82}
\mspage{II. 30}
% archive sheet number: 81
% lines 1–6
\uline{Remarque} 5.4. Dans Exp \uline{III}, nous prouverons \unsure{$=$|$\simeq$}, quand $X_K$ est une courbe \deleted{projective} \inserted{séparée} lisse sur $K$, que l'action de $I$ sur $E=H^1(X_{\bar K},\mathbb Z_\ell)$ est \unsure{pot$^{\mathrm t}$|ess$^{\mathrm t}$}-unipotente d'\uline{échelon} 2, i.e. qu'il existe un ss-groupe ouvert \deleted{\unsure{invariant}} \inserted{$U$} de $I$ tel que \deleted{\unsure{$I$}} $U$ opère \deleted{de façon} trivialement sur \deleted{\unsure{$E$}} $E/E^U$.
% lines 7–10
On voit alors facilement, en suivant les réductions 5.1, 5.2, 5.3, \deleted{qu'il en est de même} pour tout \deleted{\unsure{tt}} $X$ \deleted{normal} de type fini sur $K$, l'action de $I$ est \deleted{quasi} ess$^{\mathrm t}$ unipotente d'échelon 3 (et $=$ d'échelon 2 \deleted{\unsure{si}} si $X$ est normal).
% lines 11–18; lower half
\uline{Bibliographie}

[1] A. Grothendieck, Un th. sur les hom. de schémas abéliens, Inv. math 2, 59-78 (1966).

(2) S. Lang,

[3] J.P. Serre--J. Tate, Good reduction of abelian varieties.

[\unsure{4|3}] Shimura--Taniyama, Complex multiplication of abelian varieties, Pub. of Math. Soc. of Japan, n° 6 (1961).
